\section{现代人工智能方法综合练习}
\begin{enumerate}[leftmargin=*]
  \item \textbf{口述概念：}tensor 的 shape、dtype、device 和梯度状态分别约束什么？
  \item \textbf{口述概念：}因果卷积、RNN、self-attention 和 Transformer 如何表示历史？
  \item \textbf{口述概念：}GNN 和 DRL 分别把普通监督学习问题增加了什么数学对象？
  \item \textbf{笔试推导：}写出两层 MLP 的前向传播与平方损失梯度链，并解释残差连接为何改善梯度路径。
  \item \textbf{笔试推导：}计算核 $(1,-1)$ 对序列 $(1,3,2,5)$ 的因果卷积；再说明普通对称卷积为何可能泄漏。
  \item \textbf{笔试推导：}解释没有位置编码的 self-attention 对时间置换为何不敏感，并手算两 token softmax 权重。
  \item \textbf{笔试推导：}写出 GCN 一层传播和 Bellman 方程，分别指出图快照泄漏与环境错设发生在哪里。
  \item \textbf{数值编程：}把序列反转，比较均值、因果卷积、RNN、Attention 与 Transformer 输出，并验证 padding mask 与 causal mask 的差别。

  \item \textbf{数值编程：}用两层标量网络比较学习率 $0.1$ 与 $0.01$ 的一次更新，解释损失是否下降。
  \item \textbf{数值编程：}构造三节点供应链图，比较静态全图与逐期图快照的预测差异；不得使用预测日后新增的边。
  \item \textbf{研究判断：}审计一个把均值池化随机特征称作“深度序列模型”的研究材料。
  \item \textbf{研究判断：}比较全量微调、只训任务头和 LoRA/adapter 的证据边界，并列出 LLM 历史回测的参数记忆污染。
  \item \textbf{研究判断：}设计一个执行 RL 的最小解析或独立计算参照，并说明为什么固定历史价格回放不足以评价未执行动作。
  \item \textbf{开放项目：}从序列或文本预测、动态图预测、执行决策中选一个问题，比较简单基线与所学模型，并分析成本及一个模型假设变化的影响。
\end{enumerate}

\subsection*{两层网络的全部梯度：基础练习}
上述旧参数下若 $x=0$，哪些梯度为零？

\subsection*{记忆衰减与注意力的数值含义：基础练习}
$h_t=0.5h_{t-1}+x_t$，输入 $(1,0,0)$，求三个隐状态。

\subsection*{Bellman 递推与图上的重复平均：基础练习}
若三个节点初始特征都为 5，一次均值传播后是多少？
